\chapter{Moving Frames and Their Limits}

The preceding chapter constructed one local railway frame because the running
Alignment required a cross-section orientation. This chapter states what each
direction depends on and where the explanatory convention ends.

\section{Longitudinal Direction}

For the realized spatial trajectory \(\Gamma(s)\), the unit tangent is
\[
\mathbf t(s)=\frac{\Gamma'(s)}{\|\Gamma'(s)\|}.
\]
It follows from horizontal pose, vertical profile, metric interpretation, and
the orientation of increasing \(s\). Horizontal curvature alone does not
determine its three-dimensional direction.

\section{Uncanted Cross-Section}

The realization supplies a usable vertical reference
\(\mathbf z_{\mathrm{ref}}(s)\): locally this may be the vertical of a planar
engineering realization; in an Earth realization it may involve surface
normal and gravity conventions. Projecting that reference orthogonally to
\(\mathbf t\), normalizing it, and completing a right-handed triad gives one
possible \(F_0=(\mathbf t,\mathbf l_0,\mathbf z_0)\). A continuous-transport
rule is required where the projection becomes ill-conditioned.

This construction is deliberately qualified. The Thesis does not choose one
universal railway frame for every realization. What must remain explicit is:
\begin{itemize}
\item the longitudinal direction follows the realized trajectory and
  parameter orientation;
\item the uncanted lateral--vertical pair follows a stated vertical and
  transport convention;
\item \(u(s)\) or \(\phi(s)\) rotates that pair about \(\mathbf t\);
\item coordinate components follow the selected representation.
\end{itemize}

\section{Why a Pure Frenet Frame Is Insufficient}

A Frenet normal derived only from curvature is undefined or numerically
unstable as \(\kappa\rightarrow0\). The running Alignment contains an initial
straight, so a railway frame must carry orientation continuously across that
segment. This is not a new architectural object: it is an applicability
condition of the frame-evaluation operation.

\section{Degrees of Determination}

At a fixed \(s\), Alignment identity and the ordered law are preserved.
The starting pose and metric realization fix the planar placement. The
profile fixes elevation and slope where defined. The cant law fixes roll only
under a gauge and sign convention. A residual yaw or cross-section convention
may remain free until the selected frame rule resolves it. Vehicle body
orientation remains free until a vehicle model and operating state are
provided.

\section{World--Track and Vehicle-Space Use}

World--track projection consumes the qualified frame: \(\mathbf t\) defines
longitudinal direction, \(\mathbf l_\phi\) lateral offset, and
\(\mathbf z_\phi\) cross-section height. A vehicle-space interpretation may
then consume this track frame, but it also requires vehicle-specific data. The
frame preserves the Alignment reference and provenance of its realization; it
cannot decide whether a measured point belongs to the intended object or
whether a clearance result is acceptable.

\begin{summary}
A moving railway frame is an evaluated relation between intrinsic geometry,
vertical realization, cant, and convention. Its components are useful only
with those dependencies attached.
\end{summary}
